\chapter*{Đề cương chi tiết}
\label{decuong}

\begin{itemize}
    \item \textbf{Chương 1 - Giới thiệu}: Nêu lý do chọn đề tài, giới thiệu bài toán phân đoạn ảnh y sinh, động lực nghiên cứu, mục tiêu, đối tượng, phạm vi nghiên cứu (giả lập lượng tử), ý nghĩa thực tiễn và các đóng góp chính của khóa luận.
    \item \textbf{Chương 2 - Công trình liên quan}: Tóm tắt các nghiên cứu trước đây về Deep Learning (đặc biệt là U-Net) cho xử lý ảnh, các công trình Quantum Machine Learning cho ảnh y khoa, và cơ chế Mixture of Experts trong học sâu.
    \item \textbf{Chương 3 - Cơ sở lý thuyết}: Trình bày lý thuyết nền tảng về ảnh y sinh, kiến trúc mạng U-Net, không gian Hilbert, qubit, cổng lượng tử, mã hóa dữ liệu lượng tử (Amplitude Embedding) và nguyên lý hoạt động của Mixture of Experts.
    \item \textbf{Chương 4 - Phương pháp đề xuất}: Trình bày chi tiết kiến trúc lai QuNet, bao gồm các biến thể QuNet 12-1 và 12-2. Đi sâu vào mô tả khối Quantum Mixture of Experts (QuantumMoE), mạng định tuyến (Router), Quantum Skip Connection (QSkip) và Quantum Activation (QAct). Nêu phương pháp đánh giá so sánh.
    \item \textbf{Chương 5 - Thực nghiệm}: Mô tả chi tiết cấu hình thực nghiệm, cài đặt môi trường mô phỏng (PennyLane), các tập dữ liệu sử dụng (PH2, ISBI 2012, ISIC 2018). Trình bày kết quả đánh giá (IoU, Dice) và các phân tích bóc tách (Ablation study) cho từng thành phần lượng tử.
    \item \textbf{Chương 6 - Kết luận và hướng phát triển}: Tóm tắt kết quả, đánh giá mức độ hoàn thành mục tiêu. Chỉ ra những hạn chế hiện tại (về tài nguyên giả lập) và đề xuất các cải tiến, ứng dụng lên máy tính lượng tử thực tế trong tương lai.
\end{itemize}
